\chapter*{Lời Mở Đầu}
\addcontentsline{toc}{chapter}{Lời Mở Đầu}
\markboth{Lời Mở Đầu}{Lời Mở Đầu}

\begin{truyen}{Lời Mở Đầu}{Opening Words}
\chuhoa{K}{hông có điều gì là hoàn toàn tốt hay hoàn toàn xấu — điều quan trọng là ta nhìn nó từ góc nào.}
\textit{(Nothing is entirely good or entirely bad --- what matters is the angle from which we look at it.)}

Quyển năm của bộ sách đặt ra câu hỏi về hai mặt của mọi thứ. Mỗi câu chuyện được cặp đôi với một câu chuyện đối lập — để người đọc thấy rằng sự thật thường nằm ở giữa, không phải ở một phía.
\textit{(Volume five raises the question of the two sides of everything. Each story is paired with an opposing story --- so readers see that the truth often lies in the middle, not on one side.)}

Tư duy biện chứng không phải là kỹ năng của người lớn. Ngay từ trong trường học, học sinh đã cần học nhìn vấn đề từ nhiều chiều.
\textit{(Dialectical thinking is not just an adult skill. Even in school, students need to learn to view problems from multiple angles.)}

Quyển này mới lạ hơn những quyển trước — không phải chỉ kể chuyện, mà còn mời người đọc tranh luận với chính mình.
\textit{(This volume is more unusual than the previous ones --- it does not just tell stories, but invites readers to debate with themselves.)}
\end{truyen}

\begin{baihoc}
Trong mỗi điều xấu đều ẩn mầm tốt. Trong mỗi điều tốt đều tiềm ẩn mầm xấu. Người trí biết nhìn cả hai.
\textit{(In every bad thing lies the seed of good. In every good thing lies the seed of bad. The wise person sees both.)}
\end{baihoc}

\begin{tcolorbox}[colback=truyenblue!6,colframe=truyenblue,coltitle=white,fonttitle=\bfseries,title={Easy English},arc=4pt,boxrule=1pt,breakable]
\textbf{Short English to remember:}

Every coin has two sides.

Seeing both sides is the beginning of wisdom.
\end{tcolorbox}

\begin{tcolorbox}[colback=truyengold!10,colframe=truyengold!70!black,coltitle=black,fonttitle=\bfseries,title={Tự Học Nhanh},arc=4pt,boxrule=1pt,breakable]
1. Hãy chọn một sự kiện trong cuộc sống và thử nhìn nó từ hai góc hoàn toàn đối lập. \textit{Choose one event in your life and try to view it from two completely opposite angles.}

2. Câu chuyện nào trong quyển này thay đổi cách em nhìn một điều vốn em nghĩ là đơn giản? \textit{Which story in this volume changed how you see something you thought was simple?}

3. Khi nào tư duy một chiều có thể gây hại? \textit{When can one-sided thinking cause harm?}
\end{tcolorbox}

\ngancach
